\mxunit{automatismes-ch09}{Croisement de deux variables qualitatives}

\begin{fluencysheet}{jeu.~11~mars~2027}{Croisement de deux variables qualitatives}
\mxfq{Échauffement · Flash}{PRO}{Exprimer $\dfrac{24}{60}$ en pourcentage.}{$\pc{40}$.}{}
\mxfq{Cœur · Flash}{LEC}{Dans une classe de $30$ élèves, $18$ sont des filles. Quelle est la fréquence des filles\,? des garçons\,?}{$\dfrac{18}{30}=\pc{60}$\,; $\pc{40}$.}{}
\mxfq{Cœur + · Classer}{STA}{Classer chaque variable~: couleur des yeux\,; mention au baccalauréat (AB, B, TB)\,; taille en cm.}{Qualitative nominale\,; qualitative ordinale (ordre naturel)\,; quantitative.}{}
\mxfq{Retour · Problème}{MOD}{Un rectangle de périmètre $20$ cm a pour longueur $x$. Exprimer son aire $A(x)$\,; calculer $A(3)$ et $A(5)$.}{Largeur $10-x$\,; $A(x)=x(10-x)$\,; $A(3)=21$, $A(5)=25$ cm$^2$.}{}
\mxfq{Fil rouge · Flash}{SOL}{Calculer le volume exact d'un cylindre de rayon $2$ cm et de hauteur $5$ cm.}{$V=\pi\times4\times5=20\pi$ cm$^3$.}{}
\end{fluencysheet}

\begin{fluencysheet}{lun.~15~mars~2027}{Variables qualitatives\,; tableau croisé d'effectifs}
\mxfq{Échauffement · Flash}{CAL}{Compléter~: $12+\ldots=30$ et $22+\ldots=40$.}{$18$\,; $18$.}{}
\mxfq{Cœur · Tableau}{LEC}{Total $40$ élèves. Filles~: $22$\,; demi-pensionnaires (DP)~: $25$\,; filles DP~: $15$. Compléter le tableau croisé.\par\smallskip\mxarraystretch{1.15}\begin{tabular}{|c|c|c|c|}\hline
 & DP & Externes & Total \\ \hline
Filles & 15 & \ldots & 22 \\ \hline
Garçons & \ldots & \ldots & \ldots \\ \hline
Total & 25 & \ldots & 40 \\ \hline
\end{tabular}\par }{Filles externes $7$\,; garçons DP $10$\,; garçons externes $8$\,; total garçons $18$\,; total externes $15$ (contrôle~: $7+8=15$).}{}
\mxfq{Cœur + · Vrai/faux}{LEC}{Vrai ou faux\,? Justifier. «~Niveau de satisfaction (faible, moyen, élevé)~» est une variable qualitative nominale.}{Faux. Il existe un ordre naturel~: c'est une variable qualitative ordinale.}{}
\mxfq{Retour · Problème}{CMP}{Réduction A~: $-\pc{30}$ sur $\num{80000}$ Ar. Réduction B~: $-\pc{20}$ puis $-\pc{10}$. Laquelle est la plus avantageuse\,?}{A~: $\num{56000}$ Ar\,; B~: $\num{80000}\times\num{0,8}\times\num{0,9}=\num{57600}$ Ar. A est plus avantageuse.}{}
\mxfq{Fil rouge · Méthode}{DEM}{On sait que $AB^2+AC^2=BC^2$. Quelle conclusion tire-t-on\,? Quelle propriété\,?}{$ABC$ est rectangle en $A$~: réciproque du théorème de Pythagore.}{}
\end{fluencysheet}

\begin{fluencysheet}{mar.~16~mars~2027}{Fréquences marginales}
\mxfq{Échauffement · Flash}{PRO}{Exprimer $\dfrac{25}{40}$ en pourcentage.}{$\pc{62,5}$.}{}
\mxfq{Cœur · Méthode}{LEC}{Avec le tableau de la fiche 50~: calculer la fréquence marginale des filles et celle des demi-pensionnaires.}{Filles~: $\dfrac{22}{40}=\pc{55}$\,; DP~: $\dfrac{25}{40}=\pc{62,5}$ (effectifs marginaux sur effectif total).}{Fréquence marginale = total de la ligne (ou colonne) / effectif total.}
\mxfq{Cœur + · Flash}{LEC}{La fréquence marginale des filles est $\pc{55}$. Quelle est celle des garçons\,?}{$\pc{45}$~: les fréquences marginales d'une variable somment à $\pc{100}$.}{}
\mxfq{Retour · Flash}{SUB}{Calculer $(2a-b)^2$ pour $a=3$ et $b=-2$.}{$(6+2)^2=64$.}{}
\mxfq{Fil rouge · Flash}{REP}{$A\pt{-1;2}$ et $B\pt{3;2}$. Que peut-on dire de la droite $(AB)$\,? Donner $AB$.}{Horizontale (même ordonnée), parallèle à l'axe des abscisses\,; $AB=3-(-1)=4$.}{}
\end{fluencysheet}

\begin{fluencysheet}{jeu.~18~mars~2027}{Fréquences conditionnelles}
\mxfq{Échauffement · Flash}{PRO}{Donner $\dfrac{15}{22}$ en pourcentage arrondi au dixième.}{$\pc{68,2}$.}{}
\mxfq{Cœur · Méthode}{LEC}{Avec le même tableau~: parmi les filles, quelle fréquence de DP\,? Parmi les garçons\,?}{Filles~: $\dfrac{15}{22}\approx\pc{68,2}$\,; garçons~: $\dfrac{10}{18}\approx\pc{55,6}$ (on divise par le total de la sous-population).}{}
\mxfq{Cœur + · Distinguer}{LEC}{Distinguer «~la fréquence des filles parmi les DP~» et «~la fréquence des DP parmi les filles~».}{$\dfrac{15}{25}=\pc{60}$ et $\dfrac{15}{22}\approx\pc{68,2}$~: le total de référence n'est pas le même.}{}
\mxfq{Retour · Programme de construction}{CON}{On trace $[AB]$ de $6$ cm, puis les cercles de centres $A$ et $B$ et de rayon $4$ cm, qui se coupent en $C$ et $D$. Nature de $ACBD$\,?}{Un losange~: $AC=CB=BD=DA=4$ cm (quatre côtés égaux).}{}
\mxfq{Fil rouge · Flash}{DIV}{$91$ est-il premier\,?}{Non~: $91=7\times13$.}{Piège~: il n'est divisible ni par $2$, ni par $3$, ni par $5$.}
\end{fluencysheet}

\begin{fluencysheet}{lun.~22~mars~2027}{Compléter un tableau croisé}
\mxfq{Échauffement · Flash}{PRO}{Calculer $\pc{35}$ de $80$.}{$28$.}{}
\mxfq{Cœur · Méthode}{LEC}{$200$ clients~: $\pc{60}$ de femmes\,; $\pc{25}$ des femmes et $\pc{50}$ des hommes achètent en ligne. Compléter le tableau croisé.}{Femmes $120$ dont $30$ en ligne\,; hommes $80$ dont $40$ en ligne\,; en ligne au total~: $70$.}{On passe des fréquences conditionnelles aux effectifs en multipliant par l'effectif de la sous-population.}
\mxfq{Cœur + · Contrôle}{LEC}{Un tableau annonce «~total des lignes~: $45+55=100$~» mais «~total des colonnes~: $60+30=90$~». Que conclure\,?}{Le tableau est incohérent~: le total général doit être le même par lignes et par colonnes\,; il y a une erreur de saisie ou de calcul.}{}
\mxfq{Retour · Vrai/faux}{PRB}{Vrai ou faux\,? Justifier. Après trois «~pile~» consécutifs, la probabilité d'obtenir «~face~» au lancer suivant est supérieure à $\dfrac12$.}{Faux. Les lancers sont indépendants~: la probabilité reste $\dfrac12$.}{}
\mxfq{Fil rouge · Flash}{GEO}{$ABC$ rectangle en $A$, $BC=10$ cm, $\widehat{ABC}=30\degre$. Calculer $AC$.}{$AC=BC\times\sin30\degre=10\times\num{0,5}=5$ cm.}{Utilise le sinus~: côté opposé / hypoténuse.}
\end{fluencysheet}

\begin{fluencysheet}{mar.~23~mars~2027}{Données réelles~: traitement d'un fichier et filtrage (ET, OU, NON)}
\mxfq{Échauffement · Flash}{LEC}{Dans un fichier de $\num{1000}$ lignes, $400$ vérifient $A$, $300$ vérifient $B$ et $100$ vérifient $A$ et $B$. Combien vérifient «~$A$ ET $B$~»\,? «~$A$ OU $B$~»\,?}{$100$\,; $400+300-100=600$ (on ne compte pas deux fois l'intersection).}{}
\mxfq{Cœur · Méthode}{LEC}{Avec les mêmes données~: combien vérifient «~NON $A$~»\,? «~$A$ ET NON $B$~»\,?}{$\num{1000}-400=600$\,; $400-100=300$.}{}
\mxfq{Cœur + · Lecture de code}{ALG}{Quelle est la valeur finale de \code{compte}\,?\mxpseudo{compte = 0\\ for age in [15, 16, 17, 16, 18]:\\ \mxtab if age \textgreater{}= 16 and age \textless{} 18:\\ \mxtab\mxtab compte = compte + 1}}{$3$ (les âges $16$, $17$, $16$).}{}
\mxfq{Retour · Flash}{NUM}{Donner l'écriture décimale de $\dfrac78$ puis la fraction irréductible égale à $\num{0,12}$.}{$\num{0,875}$\,; $\dfrac{12}{100}=\dfrac{3}{25}$.}{}
\mxfq{Fil rouge · Formule → tableau}{REG}{$f(x)=x^2-2x$. Compléter un tableau pour $x\in\ens{0;1;2;3}$.}{$0\ ;\ -1\ ;\ 0\ ;\ 3$.}{}
\end{fluencysheet}

\begin{fluencysheet}{jeu.~25~mars~2027}{Lecture critique de formulations issues des médias}
\mxfq{Échauffement · Flash}{PRO}{Une hausse de $\pc{10}$ sur $40$~: quelle valeur finale\,?}{$40\times\num{1,10}=44$.}{}
\mxfq{Cœur · Esprit critique}{LEC}{Un article affirme~: «~$\pc{60}$ des accidents impliquent des conducteurs sobres~». Peut-on en conclure que conduire sobre est dangereux\,?}{Non~: il manque la population de référence. Si la grande majorité des conducteurs est sobre, il est normal qu'ils soient majoritaires parmi les accidentés. Il faudrait comparer des fréquences conditionnelles.}{}
\mxfq{Cœur + · Vrai/faux}{LEC}{Vrai ou faux\,? Justifier. «~Deux fumeurs sur trois sont des hommes~» signifie que deux hommes sur trois fument.}{Faux. On confond $P_{\text{fumeur}}(\text{homme})$ et $P_{\text{homme}}(\text{fumeur})$~: les populations de référence diffèrent.}{}
\mxfq{Retour · Vrai/faux}{STA}{Vrai ou faux\,? Justifier. La médiane d'une série est toujours égale à sa moyenne.}{Faux. Contre-exemple~: $1,\,2,\,9$~: moyenne $4$, médiane $2$.}{}
\mxfq{Fil rouge · Lecture de code}{ALG}{Quelle est la valeur de \code{s} à la fin\,?\mxpseudo{s = 0\\ for i in range(1, 5):\\ \mxtab s = s + i}}{$s=1+2+3+4=10$.}{Piège~: \code{range(1, 5)} s'arrête à $4$\,; $5$ est exclu.}
\end{fluencysheet}
